\section{TinyML Development and Deployment Workflow}
\label{sec:ch3-tinyml-workflow}

The TinyML workflow used logged current-waveform features to train two embedded
models: a load-context model and a series arc-fault model. The context model
estimated the appliance family during the early stable portion of a load run.
Its load-family output was then used as an input to the arc model, where it was
encoded as context fields. The arc model combined selected arc-sensitive
features with that context output to classify whether a frame was consistent
with a series arc fault.
Random Forest classification was selected because it can model nonlinear
feature interactions while remaining exportable as a compact tree ensemble for
microcontroller execution [68], [69].
\begin{figure}[H]
  \centering
  \resizebox{0.75\textwidth}{!}{%
  \begin{tikzpicture}[node distance=0.30in and 0.42in]
    \node[thesis process,text width=1.38in] (bench) {Bench tests\\normal, start, arc};
    \node[thesis process,text width=1.38in,right=of bench] (log) {Firmware CSV\\feature logs};
    \node[thesis process,text width=1.38in,right=of log] (label) {Manual review\\labels and notes};

    \node[thesis process,text width=1.38in,below=0.36in of label] (prepare) {Python data\\cleaning and splits};
    \node[thesis process,text width=1.38in,left=of prepare] (context) {Context model\\training};
    \node[thesis process,text width=1.38in,left=of context] (arc) {Random Forest\\arc training};

    \node[thesis process,text width=1.38in,below=0.36in of arc] (metrics) {Metrics\\confusion reports};
    \node[thesis process,text width=1.38in,right=of metrics] (export) {Exported C\\model headers};
    \node[thesis process,text width=1.38in,right=of export] (deploy) {Firmware\\integration and tests};

    \draw[thesis arrow] (bench) -- (log);
    \draw[thesis arrow] (log) -- (label);
    \draw[thesis arrow] (label) -- (prepare);
    \draw[thesis arrow] (prepare) -- (context);
    \draw[thesis arrow] (context) -- (arc);
    \draw[thesis arrow] (arc) -- (metrics);
    \draw[thesis arrow] (metrics) -- (export);
    \draw[thesis arrow] (export) -- (deploy);
  \end{tikzpicture}}
  \caption{Data logging, labeling, and TinyML deployment workflow.}
  \label{fig:ch3-data-labeling-deploy}
\end{figure}

\textit{Figure~\ref{fig:ch3-data-labeling-deploy}} shows that the training workflow was
not performed directly on the microcontroller. The firmware generated CSV
records from bench tests, the records were reviewed and labeled, and Python
scripts then cleaned the rows, formed training and test splits, trained the
context and arc models, computed confusion reports, and exported C headers.
The final embedded test was therefore a deployment check of the exported
headers inside the firmware, not a new training step on the ESP32-S3. This
follows the usual TinyML pattern of offline training followed by
resource-constrained on-device execution [43], [73].
\begin{figure}[H]
  \centering
  \resizebox{0.56\textwidth}{!}{%
  \begin{tikzpicture}[node distance=0.24in and 0.60in]
    \node[thesis process,text width=1.75in] (frame) {Current waveform\\512-sample frame};
    \node[thesis process,text width=1.75in,below=of frame] (features) {Computed/logged\\20 waveform fields};
    \node[thesis process,text width=1.75in,below=of features] (context) {Context model\\6 selected inputs};
    \node[thesis process,text width=1.75in,right=of context] (family) {Load-family\\code and confidence};
    \node[thesis process,text width=1.75in,below=of context] (arcvec) {Arc input vector\\6 waveform + 7 context};
    \node[thesis process,text width=1.75in,below=of arcvec] (rf) {Random Forest\\arc probability};
    \node[thesis process,text width=1.75in,below=of rf] (persist) {Persistence\\counter and latch};
    \node[thesis process,text width=1.75in,below=of persist] (relay) {Local relay,\\alerts, and logs};

    \node[thesis process,text width=1.75in,right=of features] (pwa) {Firebase/PWA\\10 live feature cards\\logging and control};

    \draw[thesis arrow] (frame) -- (features);
    \draw[thesis arrow] (features) -- (context);
    \draw[thesis arrow] (context) -- (arcvec);
    \draw[thesis arrow] (context) -- (family);
    \draw[thesis arrow] (family) |- (arcvec);
    \draw[thesis arrow] (arcvec) -- (rf);
    \draw[thesis arrow] (rf) -- (persist);
    \draw[thesis arrow] (persist) -- (relay);
    \draw[thesis arrow] (features) -- (pwa);
  \end{tikzpicture}}
  \caption{TinyML processing sequence from feature logging to embedded inference.}
  \label{fig:ch3-tinyml-sequence}
\end{figure}

\textit{Figure~\ref{fig:ch3-tinyml-sequence}} shows that the firmware computes feature
fields first, then updates context, then evaluates the arc classifier. This
order is important because the deployed arc header includes both waveform
features and context fields. The load-family result is therefore not only a
display label; when confidence is sufficient, it becomes part of the input
context supplied to the arc model. The flowchart also separates three feature
sets: the candidate fields computed and logged by the firmware, the six base
waveform features selected for the deployed arc model, and the 10 live
feature cards shown by the PWA. A broader 14-feature order remains in the PWA
metadata and session viewer for historical review and CSV inspection.

Arc labels were assigned from logged waveform sessions. The labeling interface
is shown in \textit{Figure~\ref{fig:ch3-arc-labeling}}.
\begin{figure}[H]
  \centering
  \includegraphics[width=0.99\textwidth]{ch3/figure_3.28.png}
  \caption{Arc labeling using logged waveform data.}
  \label{fig:ch3-arc-labeling}
\end{figure}

\textit{Figure~\ref{fig:ch3-arc-labeling}} shows the interface used to mark arc
events from the logged waveform sessions. The labels were based on the
observed test action and on the expected behavior of series arcs during
contact separation and restriking. Because the marking was done manually
after the test, the selected points are not exact arc-onset locations, but
time-aligned approximations near the RMS waveform changes associated with
the disconnection and recontact events.

To account for this labeling uncertainty, the training set did not treat
only the marked row as the arc event. Neighboring rows around each arc marker
were also given soft arc weights so that brief timing delays in the manual
marking would not force nearby arc-related frames to be learned as normal
behavior. This is presented in \textit{Appendix F} and is also shown in 
\textit{Figure~\ref{fig:ch3-python-trainer}}.
\begin{figure}[H]
  \centering
  \includegraphics[width=0.99\textwidth]{ch3/figure_3.29.png}
  \caption{Local Python TinyML trainer.}
  \label{fig:ch3-python-trainer}
\end{figure}

\textit{Figure~\ref{fig:ch3-python-trainer}} documents the offline training environment
used to prepare the exported embedded model files. The workflow used separate
data preparation, context-model training, Random Forest training, shared
TinyML utility, and generation-path routines as the source of truth for
feature selection, data splits, hyperparameter search, metric computation,
and embedded export. The dataset-preparation, training-report, and exported
embedded-model summaries are provided in
\textit{Appendix~\ref{app:f-tinyml-training-summary}}.

\subsection{Context Model}
\label{sec:ch3-context-model}

The load context model is a centroid classifier, not a
Random Forest. The context model defines six input features:
$\Delta I_{\mathrm{rms}}$, half-cycle asymmetry, midband residual
ratio, current THD, mid/high-frequency spectral flux, and high-frequency
energy delta. These features were used because they describe the waveform
behavior of the connected load rather than a single fault event.
For example, different load families may exhibit different harmonic content,
asymmetry, and current-change patterns even during normal operation.

In this study, the context model was used to provide load-family information
to the arc classifier, but its experimental coverage was still limited by the
available tested loads. As stated in the limitations of the study, the actual
collected and tested load signatures were limited to resistive loads,
inductive loads, and brushed motor loads.

During runtime, the firmware averages usable feature frames over the
acquisition interval, standardizes each averaged feature, computes the
squared Euclidean distance to each family centroid, and applies a softmax to
the negative distances:
\begin{equation}
  p_c =
  \frac{\exp(-d_c^2)}
       {\sum_{j=1}^{C}\exp(-d_j^2)},
  \qquad
  d_c^2 = \sum_{i=1}^{m}
  \left(\frac{x_i-\mu_i}{\sigma_i}-q_{c,i}\right)^2 .
  \label{eq:context-prototype}
\end{equation}
In \eqref{eq:context-prototype}, $p_c$ is the context confidence for family
$c$, $x_i$ is the averaged runtime feature, $\mu_i$ and $\sigma_i$ are the
training-set mean and standard deviation exported in the header, $q_{c,i}$ is
the standardized centroid coordinate for family $c$, $m=6$ is the context
input dimension, and $C=6$ is the number of load families. A family is accepted
only when the best confidence is at least 0.45; otherwise, the context remains
unknown. This threshold prevents weak or ambiguous feature patterns from being
forced into a load family.

The context acquisition constants in the firmware require a load to remain
active before the context is latched. A provisional context may be used after
the minimum acquisition period, and the final context is latched after the
configured 5s acquisition window when confidence is sufficient. This prevents
brief startup artifacts from immediately forcing a load family into the arc
classifier. Once the context is latched, it remains stable for the continuous
load run, so the arc model receives a consistent family condition during
normal monitoring. The context is returned to unknown only after the load has
been inactive long enough to indicate that the previous operating condition
has ended.

\subsection{Arc Model}
\label{sec:ch3-arc-model}

The deployed series-arc classifier is the Random Forest used by the embedded
arc model. The deployed Random Forest arc model defines 13 inputs: six
base waveform features, six one-hot context-family fields, and one context
confidence field. The base features are maximum low-current run, zero-dwell
ratio, low-current ratio, absolute RMS-current z-score relative to baseline,
midband residual ratio, and mid/high-frequency spectral flux. These were
selected because series arcs can produce current interruption, restrike spikes,
cycle-to-cycle waveform disturbance, and broadband spectral changes [34],
[35], [50], [53]. Since the reviewed literature also shows that series arcs
are load-dependent and difficult to detect using current magnitude alone, the
arc decision was implemented as a multi-feature classifier rather than as a
single-threshold current decision [33], [34], [36].

The first three base features describe current dropouts and low-current
intervals that may occur when a contact separates. The RMS-current z-score
compares the present current behavior with the recent baseline, allowing the
model to detect abrupt deviations from the expected operating level. The
midband residual ratio and mid/high-frequency spectral flux describe waveform
disturbance and spectral changes that may occur during restriking. Together,
these features allow the model to observe both the time-domain interruption
and the frequency-domain disturbance caused by a possible arc event. This
follows the reviewed arc-fault studies, where zero-current behavior,
current-fluctuation features, high-frequency noise, and spectral changes were
used as indicators of series arcing [34]-[39], [49]-[53].

The context-family inputs allow the arc classifier to interpret the same
waveform features differently depending on the connected load type. This is
important because some normal loads, such as motors or phase-controlled
devices, can naturally produce distorted or non-sinusoidal current waveforms.
The context confidence is included so that the model can also account for how
reliable the assigned load family is. When the context is unknown or weak, the
arc decision is treated more conservatively to reduce false operation during
unstable startup or ambiguous load behavior. This follows the reviewed NILM
and load-family literature, where load signatures were shown to affect how
waveform features should be interpreted [8], [40]-[42].

Random Forest was used because the reviewed TinyML literature supports compact
feature-based models that can be deployed on constrained embedded hardware
when the feature space and runtime are carefully designed [32], [43]. In this
study, the Random Forest therefore serves as the embedded decision layer that
combines the selected arc-related features and the estimated load context.
\begin{table}[H]
  \centering
  \caption{Deployed Random Forest arc-model settings}
  \label{tab:ch3-arc-model-export}
  \begin{tabular}{p{0.36\textwidth}p{0.22\textwidth}p{0.32\textwidth}}
    \toprule
    \textbf{Item} & \textbf{Deployed value} & \textbf{Source or role} \\
    \midrule
    Model type & Random Forest & Exported tree ensemble from the local training workflow \\
    Input dimension & 13 & Six waveform features plus seven context fields \\
    Base feature count & 6 & Arc-sensitive features used before context fields \\
    Default arc threshold & 0.4800 & Probability threshold for supported context families \\
    Context confidence minimum & 0.4500 & Minimum accepted context confidence \\
    Unknown-context threshold & 0.7800 & More conservative threshold when context is unknown \\
    Unknown minimum feature votes & 3 & Additional runtime safeguard for unknown-context decisions \\
    \bottomrule
  \end{tabular}
\end{table}

\textit{Table~\ref{tab:ch3-arc-model-export}} shows that the deployed inference was
context-aware. The model used lower thresholds for selected known families and
a higher threshold for unknown contexts, reducing nuisance trips when the load
family could not be established confidently. The firmware then passed the arc
classification to the protection manager, where a leaky counter required
persistence before relay disconnection.

\subsection{Feature Selection for the Deployed Arc Model}
\label{sec:ch3-arc-feature-selection}

The exported firmware header is the source of truth for the deployed feature
selection. It defines 13 arc-model inputs: six waveform features and seven
context fields. Therefore, this paper treats the context one-hot fields
and context confidence as model inputs, but not as additional waveform
features. The selected waveform features are maximum low-current run,
zero-dwell ratio, low-current ratio, absolute RMS-current z-score relative to
baseline, midband residual ratio, and mid/high-frequency spectral flux. The
full computed/logged feature contract is given in
\textit{Section~\ref{sec:ch3-waveform-features}}.

\subsection{Training Split and Model Evaluation}
\label{sec:ch3-training-split}

The Random Forest report, as documented in \textit{Appendix F}, shows that the final arc dataset was split into
35,512 training rows, 9,434 validation rows, and 10,268 test rows. The final
training report used 5 cross-validation splits during search and reported the
held-out test metrics after model selection. The best exported hyperparameters
were 160 trees, maximum tree depth of 16, minimum leaf size of 4, minimum split
size of 12, maximum feature fraction of 0.35, bootstrap sampling enabled, and
entropy criterion. Train/test splitting, cross-validation, and confusion-matrix
metrics are standard model-evaluation practices for supervised classification
[69], [74].

The held-out test-set confusion matrix in the exported report contained 10,173
true negatives, 5 false positives, 25 false negatives, and 65 true positives.
The corresponding frame-level accuracy, precision, recall, and F1-score were
99.7078\%, 92.8571\%, 72.2222\%, and 81.25\%, respectively. Event-level
evaluation was also reported because arc faults occur in grouped episodes
rather than as isolated independent frames; the event-level recall was 84.0\%,
with one false-alarm event in the test sessions.
